\scanpage{93}

Soit $X_{\alpha\beta}=X_\alpha\cap f^{-1}(Y_\beta)$, $Y_{\deleted{\unsure{\alpha\beta}}\beta\alpha}=Y_\beta\cap g^{-1}(X_\alpha)$, donc $X_{\alpha\beta}$ et $Y_{\beta\alpha}$ sont des ens \uline{constructibles}, qui se correspondent par $f,g$.

\deleted{Donc \inserted{\unsure{disons}} $X_{\alpha\beta}$ est une union finie $\bigcup_\lambda X_{\alpha\beta\lambda}$, où $X_{\alpha\beta\lambda}$ est un sous-schéma affine constructible de $X_{\alpha\beta}$ \inserted{\unsure{donc de} $X_\alpha$}. Alors $g^{-1}(X_{\alpha\beta\lambda})=g_\beta^{-1}(X_{\alpha\beta\lambda})$ est \deleted{\unsure{une partie|un sous schéma}} un sous-schéma de $Y_\beta$ constructible.}

\diagram{\begin{tikzcd}[ampersand replacement=\&,row sep=2pt,column sep=small]
X\&Y\\
{}\&Y_\beta\arrow[ul]
\end{tikzcd}}

D'ailleurs $X_{\alpha\beta}=f_\alpha^{-1}(Y_\beta)$ est un \deleted{\unsure{sous-schéma de|ens. schématique de}} loc. fermé dans $X_\alpha$, donc dans $X$. et $f_{\alpha\beta}:X_{\alpha\beta}\to Y$ \deleted{\unsure{est}}est un morphisme \inserted{de schémas} (car induit par $f_\alpha:X_\alpha\to Y$]. \deleted{\unsure{comme il en est}}\inserted{De $=$,} les $Y_{\alpha\beta}$ sont constructibles et loc. fermés, et \inserted{$g_{\alpha\beta}:Y_{\alpha\beta}\to X$ est un morphisme de schém.} \deleted{\unsure{comme}}\inserted{Enfin} $X_{\alpha\beta}\unsure{\text{---}|\to}Y_{\unsure{\alpha|\delta|\beta}}$ se factorise par un \uline{morph.} $X_{\alpha\beta}\to Y_{\alpha\beta}$ [car il le fait \unsure{ens|ainsi}], et de $=$ on a un \uline{morphisme} $Y_{\alpha\beta}\to X_{\alpha\beta}$. Ces derniers sont donc inverses l'un de l'autre.

\uline{Cor.} 0 Si $X$ noethérien $Y$ l'est
% slanted
\margin{\uline{Notion} d'ens schématique \uline{noethérien} [défini par un $X$ noethérien, \unsure{ou|qui} \illegible{1} \unsure{qu'on peut} \unsure{supposer} \unsure{affine}]}

\uline{Cor.} 1 Tout ens schématique est isomorphe à un $|X|$, avec $X$ affine.

\uline{Cor.} 2 Tout morphisme d'ens\unsure{.} schématiques est isom. à un morphisme de \deleted{ens} schémas affines. (\unsure{Considéré|Considérés} comme application \unsure{constructible|constructibles}).

Démonstration \unsure{et|tte} \unsure{réciproque|analogue|intrinsèque|logiques}.

\deleted{\uline{Cor.} 3 Un \deleted{isom}\inserted{morphisme $f:X\to Y$} de préschémas quasi compacts et quasi-séparés, qui est un isom d'ens schématiques est \uline{de type fini}}
% archive sheet number: 92
